% Generated draft: requires mathematical and source review.
\documentclass[11pt]{article}
\usepackage[margin=25mm]{geometry}
\usepackage{fontspec}
\setmainfont{Latin Modern Roman}
\usepackage{amsmath,amssymb,mathtools,unicode-math}
\setmathfont{Latin Modern Math}
\usepackage{enumitem,needspace,xurl,hyperref}
\setlength{\parindent}{0pt}
\setlength{\parskip}{5pt}
\setlength{\emergencystretch}{3em}
\newcommand{\N}{\mathbb{N}}
\newcommand{\Z}{\mathbb{Z}}
\newcommand{\Q}{\mathbb{Q}}
\newcommand{\R}{\mathbb{R}}
\newcommand{\C}{\mathbb{C}}
\newcommand{\F}{\mathbb{F}}
\newcommand{\A}{\mathbb{A}}
\newcommand{\PP}{\mathbb{P}}
\newcommand{\E}{\mathbb{E}}
\newcommand{\eps}{\varepsilon}
\newcommand{\dd}{\,\mathrm{d}}
\begin{document}
\section*{UnsolvedMath ID: 30001629}
\textbf{Source problem code:} catalogue entry
\subsection{Definitions and mathematical statement}
For a surface \(S\), a mapping class group is a group of homeomorphisms of \(S\) modulo isotopy, with the same prescribed orientation, boundary, and marked-point conditions imposed on homeomorphisms and isotopies. Write this group as \(\operatorname{Mod}(S)\) once those conditions have been specified. The supplied record does not specify the class of surfaces or these conditions.

For a countable discrete group \(G\), the \emph{Haagerup property} means that there exist a real Hilbert space \(\mathcal H\), an orthogonal representation \(\pi:G\to O(\mathcal H)\), and a map \(b:G\to\mathcal H\) satisfying
\[
 b(gh)=b(g)+\pi(g)b(h)\qquad(g,h\in G),
\]
such that
\[
 \#\{g\in G:\|b(g)\|\le R\}<\infty
 \qquad\text{for every }R\ge0.
\]
Equivalently, \(G\) admits a metrically proper affine isometric action on a real Hilbert space, given by \(g\cdot v=\pi(g)v+b(g)\).

\textbf{Source question.} To what extent do mapping class groups have the Haagerup property? Under the countable discrete interpretation just defined, the mathematical task is to determine for which intended surfaces \(S\) and mapping class group conventions such a pair \((\pi,b)\) exists for \(G=\operatorname{Mod}(S)\). This is a classification question, not an assertion that every mapping class group has the property. A fully quantified surface classification cannot be extracted from the supplied record: the surface class, boundary and marked-point conventions, and intended scope of ``to what extent'' remain unspecified. The countable discrete interpretation is explicitly a convention for this draft, not a hypothesis stated in the source.
\subsection{Short English statement}
Determine which mapping class groups admit a metrically proper affine isometric action on a Hilbert space. The supplied question leaves the surfaces and mapping class group conventions unspecified, so its precise classification scope requires review. The supplied registry labels the problem open; this draft does not independently verify that status.
\subsection{Sources}
\begin{itemize}
\item Supplied local UnsolvedMath registry record, permanent ID \(30001629\), title \emph{The Haagerup Property for Mapping Class Groups}, source problem code \texttt{catalogue entry}, in the Oberwolfach Reports Open Problems collection. The record lists proposed year \(2010\) and identifies no individual proposer.
\item Supplied source citation: \emph{Teichmüller Theory (2011)}, Oberwolfach Reports, pp.~3085--3157, DOI \texttt{10.4171/owr/2010/53}. Original OWR report: \url{https://doi.org/10.4171/owr/2010/53}. The collection describes these reports as documenting mathematical workshops and their open-problem discussions. The supplied citation uses \(2011\), while the DOI and proposed-year metadata use \(2010\); these details are retained without reconciliation.
\item Supplied literature source, labeled ``mapping-class-group coarse properties'': \url{https://numdam.org/item/10.5802/ahl.85.pdf}. No author, title, or precise supporting passage is supplied for this reference.
\item Literature assessment dated \texttt{2026-08-21}: the supplied curation labels the question open and states that recent work establishes other coarse properties while recording that the Haagerup property remains unknown. This is supplied dated curation, not an independently verified assessment. The background explicitly requests independent verification of time-sensitive status and cited claims.
\item The supplied record labels the statement \texttt{exact} and reports that its mathematical content was checked against the official Oberwolfach report, with notation, context, and question form normalized without changing the source-level claim. This verification description is supplied metadata; the sources have not been independently consulted for this draft.
\end{itemize}
\end{document}
